\documentclass[11pt,letterpaper]{article}
\usepackage[T1]{fontenc}
\usepackage[utf8]{inputenc}
\usepackage{amsmath,amsthm,mathtools}
\usepackage{newtxtext,newtxmath}
\usepackage[margin=1in,headheight=15pt]{geometry}
\usepackage{microtype}
\usepackage{booktabs,array,longtable}
\usepackage{xcolor}
\usepackage{enumitem}
\usepackage{listings}
\usepackage{fancyhdr}
\usepackage{titlesec}
\usepackage{hyperref}
\usepackage[nameinlink,noabbrev]{cleveref}
\definecolor{OliveInk}{HTML}{4E5C3D}
\definecolor{SoftOlive}{HTML}{EFF2E8}
\definecolor{Muted}{HTML}{60665D}
\definecolor{Coral}{HTML}{A85E53}
\hypersetup{colorlinks=true,linkcolor=OliveInk,citecolor=OliveInk,urlcolor=OliveInk,
 pdftitle={Single curves do not force log-concavity},
 pdfauthor={Research note prepared with ChatGPT for Vladimir Reshetnikov},
 pdfsubject={An annular counterexample, positive transfer matrices, and Fibonacci-type recurrences}}
\titleformat{\section}{\large\bfseries\color{OliveInk}}{\thesection}{0.75em}{}
\titleformat{\subsection}{\normalsize\bfseries\color{OliveInk}}{\thesubsection}{0.65em}{}
\pagestyle{fancy}\fancyhf{}
\fancyhead[L]{\small\color{Muted}Single curves and coefficient sequences}
\fancyhead[R]{\small\color{Muted}19 September 2026}
\fancyfoot[C]{\small\thepage}
\renewcommand{\headrulewidth}{0.3pt}
\setlength{\parskip}{0.25em}
\setlength{\emergencystretch}{2em}
\setlist{nosep,leftmargin=1.5em}
\numberwithin{equation}{section}
\newtheorem{theorem}{Theorem}[section]
\newtheorem{lemma}[theorem]{Lemma}
\newtheorem{proposition}[theorem]{Proposition}
\newtheorem{corollary}[theorem]{Corollary}
\theoremstyle{definition}
\newtheorem{definition}[theorem]{Definition}
\newtheorem{remark}[theorem]{Remark}
\newcommand{\Z}{\mathbb Z}
\newcommand{\N}{\mathbb N}
\newcommand{\R}{\mathbb R}
\newcommand{\Q}{\mathbb Q}
\newcommand{\Btilde}{\widetilde B}
\newcommand{\pos}[1]{[#1]_+}
\newcommand{\coeff}[2]{[q^{#1}]\,#2}
\newcommand{\Pn}{P_n^{(d)}}
\newcommand{\EE}{\mathbb E}
\newcommand{\Var}{\operatorname{Var}}
\newcommand{\diag}{\operatorname{diag}}
\lstset{basicstyle=\small\ttfamily,breaklines=true,columns=fullflexible,
 keepspaces=true,showstringspaces=false,frame=single,rulecolor=\color{OliveInk},
 backgroundcolor=\color{SoftOlive},commentstyle=\color{Muted},
 keywordstyle=\bfseries\color{OliveInk},numbers=none,tabsize=4}

\begin{document}
\begin{center}
{\small\scshape\color{OliveInk}A constructive attack on a recent combinatorial conjecture}\par
\vspace{0.7em}
{\LARGE\bfseries Single curves do not force log-concavity}\par
\vspace{0.45em}
{\large An annular counterexample, a positive transfer matrix,\\[0.15em]
and Fibonacci-type recurrences}\par
\vspace{0.9em}
{\normalsize Research note prepared with ChatGPT for Vladimir Reshetnikov}\par
{\small 19 September 2026}
\end{center}
\vspace{0.4em}
\begin{abstract}
We construct a counterexample to the single-lamination log-concavity conjecture
in Section~5 of Kang--Lee--Lim, \emph{Unimodality and Cluster Algebras from
Surfaces}, arXiv:2508.04396v3. The surface is an unpunctured annulus with one
marked point on each boundary component, and the lamination consists of one
unit-weight simple curve. A legal change of triangulation gives the extended
exchange matrix with rows $(0,-2)$, $(2,0)$, and $(-1,2)$. Three subsequent
mutations produce the coefficient polynomial
\[
1+3q^2+2q^3+4q^4+2q^5+q^6.
\]
The positive consecutive coefficients $3,2,4$ violate log-concavity and form a
strict valley, so deleting internal zeros does not remove the obstruction.
Indeed, one fixed triangulation and one fixed curve yield infinitely many
nonunimodal examples. We derive a two-parameter family, prove an equivalent
linear recurrence, give a nonnegative two-state walk model and a positive
finite coefficient formula, identify its row sums with OEIS A001519, and
establish exponential and Gaussian asymptotics. The geometric realization is
proved from standard surface-cluster results; all finite algebraic certificates
are also checked with an included dependency-free exact-arithmetic program.
The disproof targets the stated arXiv version; no claim of a completed priority
search or independent peer review is made.
\end{abstract}

\setcounter{tocdepth}{1}
\tableofcontents
\newpage

\section{The problem and the outcome}\label{sec:problem}

For a polynomial $f(q)=\sum_{i=0}^{m}a_iq^i$ with nonnegative coefficients,
\emph{log-concavity} means
\begin{equation}\label{eq:lc}
 a_i^2\geq a_{i-1}a_{i+1}\qquad(1\leq i<m).
\end{equation}
A sequence is \emph{unimodal} if it first weakly increases and then weakly
decreases. A strict local valley $a_{i-1}>a_i<a_{i+1}$ therefore disproves
unimodality, independently of what happens elsewhere.

The selected problem is Conjecture~1 of Kang--Lee--Lim~\cite{KLL}, in the
single-lamination setting of their Section~5. Their Definition~20 associates a
$c$-polynomial to a cluster variable by specializing the initial mutable
variables to $1$ and the frozen variables to $q$. They conjecture that its
coefficient sequence satisfies \eqref{eq:lc} when the lamination has one
component. Their Theorem~14 also asserts unimodality on unpunctured surfaces.
The construction below contradicts both statements at that level of generality.

\begin{theorem}[Counterexample and persistence]\label{thm:main}
There are an unpunctured marked annulus $S$, a fixed triangulation $T$, and a
fixed lamination $L$ consisting of exactly one unit-weight simple curve such
that the $c$-polynomials of a sequence of simple arcs are $P_n(q)$, where
\begin{equation}\label{eq:mainlinear}
 P_0=P_1=1,\qquad
 P_{n+1}=(1+q+q^2)P_n-qP_{n-1}\quad(n\geq1).
\end{equation}
For every $n\geq3$, the positive consecutive coefficients at degrees $2,3,4$
violate log-concavity. For every $n\geq4$, those coefficients form a strict
valley. In particular,
\begin{equation}\label{eq:P4}
 P_4(q)=1+3q^2+2q^3+4q^4+2q^5+q^6
\end{equation}
is neither log-concave nor unimodal, even after all zero coefficients are
removed.
\end{theorem}

The proof has two independent parts. The geometric part verifies that the
coefficient row really comes from \emph{one connected curve}, not an arbitrary
integral lamination. The algebraic part consists of explicit mutations and
coefficient calculations. Both are necessary: a random integer row does not
by itself certify connectedness of the corresponding lamination.

\subsection{Version and scope of the claim}

The source examined is arXiv:2508.04396v3, whose submission history gives
11 February 2026. The arXiv record lists a related publication in the
\emph{European Journal of Combinatorics}, volume~136, article~104392~\cite{KLL}.
The final publisher text was not independently inspected. Accordingly,
``counterexample to the conjecture'' throughout this note means a counterexample
to its formulation in that specific arXiv version, with its definitions and
exchange rules. The result concerns an annulus, not a polygon-only restriction,
and does not contradict the separate results about ordinary rank polynomials
of fence or loop-fence posets.

The main mathematical assertion is a proof, rather than an inference from a
search bound. The computational appendix corroborates its formulas but is not
used to replace an argument for infinitely many indices. The recurrence and
transfer-matrix techniques are elementary; we make no separate originality
claim for those methods or for the Fibonacci sequence obtained by specialization.

\section{The geometric certificate}\label{sec:geometry}

This section supplies all surface-cluster input needed later. Readers primarily
interested in enumerative combinatorics can then work entirely with the
recurrences and weighted walks.

\subsection{Mutation with one frozen variable}

A rank-two geometric seed consists of independent mutable variables $x_1,x_2$,
a frozen variable $q$, and an integer matrix
\[
 \Btilde=(b_{ij})_{1\leq i\leq3,\,1\leq j\leq2}
\]
whose top $2\times2$ submatrix is skew-symmetric. Mutation at $k\in\{1,2\}$
replaces $x_k$ by $x_k'$ according to
\begin{equation}\label{eq:exchange}
 x_kx_k'
 =\prod_{i=1}^{3}x_i^{\pos{b_{ik}}}
  +\prod_{i=1}^{3}x_i^{\pos{-b_{ik}}},
 \qquad x_3=q,
\end{equation}
where $[z]_+=\max(z,0)$. Simultaneously,
\begin{equation}\label{eq:matrixmutation}
 b'_{ij}=\begin{cases}
 -b_{ij},&i=k\text{ or }j=k,\\
 b_{ij}+[b_{ik}]_+[b_{kj}]_+
             -[-b_{ik}]_+[-b_{kj}]_+,&\text{otherwise}.
 \end{cases}
\end{equation}
These are the usual exchange and matrix-mutation rules; see
Fomin--Zelevinsky~\cite{FZ} and Fomin--Thurston~\cite{FT}.

There is just one frozen generator in our construction. Its powers are not
new variables. In particular, $q^2$ remains $q^2$ after the specialization in
Definition~20; it is not renamed $q$. This distinction matters when a single
curve has shear coordinate of absolute value greater than one.

\subsection{Two triangles on an annulus}

Let $S$ be an annulus with a marked point on each boundary circle and no
punctures. A concrete model is a square with its vertical sides identified.
The two bottom corners become one marked point and the two top corners become
the other. The diagonal and the identified vertical side give two bridging
arcs, denoted $\alpha_1$ and $\alpha_2$. Their interiors are disjoint. They cut
the annulus into two triangles, neither of which is self-folded.

With a suitable labeling, both triangles contribute the same arrow between
these two arcs. Equivalently, in the square model their cyclic side orders give
$\alpha_2$ immediately after $\alpha_1$ in the signed-adjacency convention in
each triangle. Thus the principal exchange matrix is
\begin{equation}\label{eq:B0principal}
 B(T_0)=\begin{pmatrix}0&2\\-2&0\end{pmatrix}.
\end{equation}
The entry $2$ is the sum of two contributions from different triangles. It does
not mean that an arc or a lamination has been given multiplicity two.

This annulus lies within the usual marked-surface setup: each boundary
component has a marked point, and it has no punctures or self-folded triangles.
Arcs and flips have the standard interpretation established by
Fomin--Shapiro--Thurston~\cite{FST} and Fomin--Thurston~\cite{FT}.

\subsection{One curve, not merely one coefficient row}

Take $L=L_{\alpha_1}$, the elementary lamination following $\alpha_1$ once,
with its endpoints displaced slightly along the boundary in the standard
prescribed directions. This is a single simple bridging curve of weight one.
Because its endpoints lie on different boundary components, it is not
boundary-parallel and does not bound an excluded disk.

The elementary-lamination construction gives the shear row $(1,0)$ relative
to $T_0$. This is the one-row part of the principal-coefficient realization in
Fomin--Thurston, Definition~17.2 and Proposition~17.3~\cite{FT}. We use only
this one curve, not the entire collection of elementary laminations. Hence
\begin{equation}\label{eq:B0extended}
 \Btilde_0=\begin{pmatrix}0&2\\-2&0\\1&0\end{pmatrix}.
\end{equation}

Keep $L$ fixed and flip $\alpha_1$. The shear-coordinate mutation theorem
(Fomin--Thurston, Theorem~12.6, or its tagged version Theorem~13.5) says that
changing the triangulation changes the extended matrix by
\eqref{eq:matrixmutation}~\cite{FT}. Direct calculation gives
\begin{equation}\label{eq:certificate}
 \Btilde_1=\mu_1(\Btilde_0)
 =\begin{pmatrix}0&-2\\2&0\\-1&2\end{pmatrix}.
\end{equation}
For clarity, the new bottom-right entry is
$0+[1]_+[2]_+-[-1]_+[-2]_+=2$, and the first bottom entry is negated.

\begin{lemma}[Connected geometric realization]\label{lem:realization}
The matrix \eqref{eq:certificate} is the extended exchange matrix of the
triangulation $T_1=\mu_1(T_0)$ and the same one-component, unit-weight lamination
$L=L_{\alpha_1}$.
\end{lemma}
\begin{proof}
The initial realization \eqref{eq:B0extended} uses the explicit elementary
curve just described. The flip changes only the triangulation. The fixed-curve
shear-coordinate mutation theorem gives \eqref{eq:certificate}. Neither the
number of connected components of $L$ nor its weight changes under this
operation.
\end{proof}

In particular, we have not inferred connectedness from the fact that $-1$ and
$2$ are relatively prime. Such an inference would not be a valid general
criterion. The actual curve and its legal flip history are the certificate.

\subsection{Which variables are set to one?}\label{sec:reseed}

We now choose $T_1$ as the \emph{initial triangulation for the $c$-polynomials}.
Write its independent mutable variables as $a,b$. Every $c$-polynomial below
means specialization of these variables to $a=b=1$ with $q$ retained.

This is not the same as setting the variables of $T_0$ to one and then
repeatedly resetting newly created variables. To see the distinction
explicitly, write the variables of $T_0$ as $u,b$. Its first exchange is
\[
 ua=q+b^2.
\]
Thus $a=b=1$ corresponds to $u=1+q$, not to $u=1$. The birational change of
coordinates ensures that $a,b$ remain independent over $\Q(q)$. The general
statement being tested allows a chosen initial triangulation, so choosing
$T_1$ is legitimate.

Although $L$ was conveniently described as elementary relative to $T_0$, it
is not one of the elementary coefficient curves of $T_1$: its row there is
$(-1,2)$ rather than a standard basis row. In particular this is not a
principal-coefficient rank-polynomial calculation in disguise.

\section{The counterexample in three exchanges}\label{sec:small}

Start with \eqref{eq:certificate} and the cluster $(a,b)$. Mutate in directions
$2,1,2$. Denote the successive new cluster variables by $v,w,z$. The matrices
and equations are completely explicit:
\begin{align}
 \Btilde_1&=\begin{pmatrix}0&-2\\2&0\\-1&2\end{pmatrix},
 & bv&=a^2+q^2,\label{eq:firstexchange}\\
 \Btilde_2&=\begin{pmatrix}0&2\\-2&0\\3&-2\end{pmatrix},
 & aw&=v^2+q^3,\label{eq:secondexchange}\\
 \Btilde_3&=\begin{pmatrix}0&-2\\2&0\\-3&4\end{pmatrix},
 & vz&=w^2+q^4.\label{eq:thirdexchange}
\end{align}
After the last exchange the matrix has bottom row $(5,-4)$ and principal
part $B(T_0)$. All variables created by these flips correspond to simple arcs
on the annulus. No generalized self-intersecting arc or essential closed-loop
variable is needed.

\subsection{A full Laurent certificate}

To avoid relying solely on calculations after specialization, put
$A=a^2+q^2$. The exchanges give
\begin{equation}\label{eq:laurentcertificate}
 v=\frac{A}{b},\qquad
 w=\frac{A^2+b^2q^3}{ab^2},\qquad
 z=\frac{A^3+2Ab^2q^3+b^4q^4}{a^2b^3}.
\end{equation}
For the last identity, substitute the first two expressions into
\eqref{eq:thirdexchange}. The numerator before cancellation is
\[
 (A^2+b^2q^3)^2+a^2b^4q^4
 =A\bigl(A^3+2Ab^2q^3+b^4q^4\bigr),
\]
because $q^6+a^2q^4=q^4A$. This establishes the Laurent expression directly.

Now set $a=b=1$. The successive polynomials are
\begin{align}
 P_2&=1+q^2,\label{eq:P2}\\
 P_3&=1+2q^2+q^3+q^4,\label{eq:P3}\\
 P_4&=1+3q^2+2q^3+4q^4+2q^5+q^6.\label{eq:P4again}
\end{align}
Already $P_2$ violates \eqref{eq:lc}: its coefficients are $(1,0,1)$ and
$0^2<1\cdot1$. The stronger example is $P_4$. At index $3$, its coefficients
satisfy
\begin{equation}\label{eq:positivefailure}
 (\coeff{3}{P_4})^2=2^2=4
 <3\cdot4=(\coeff{2}{P_4})(\coeff{4}{P_4}).
\end{equation}
Moreover, $3>2<4$, so the sequence is not unimodal.

\subsection{Why an internal-zero convention cannot rescue the statement}

One can additionally require the absence of internal zeros, or instead adopt
the more permissive convention of testing only the sequence of nonzero
coefficients. Neither choice matters for \eqref{eq:positivefailure}. Removing the zero at degree one
leaves
\[
 (1,3,2,4,2,1),
\]
which still has the strict valley $3,2,4$ and still violates the corresponding
log-concavity inequality. Likewise, multiplying by a monomial merely shifts
degrees and does not remove this triple.

One must not generalize the first binomial example carelessly. When $d>2$,
the sequence of $1+q^d$ can satisfy every \emph{local} inequality
\eqref{eq:lc} if internal zeros are allowed, because no zero has two positive
neighbors. This is why the positive-consecutive-coefficient obstruction is
preferable to an argument based only on a support gap.

\section{A two-parameter annular family}\label{sec:family}

The three-exchange calculation is not isolated. The same curve yields
arbitrarily large shear coordinates as the triangulation changes.

\begin{lemma}[Alternating shear-coordinate orbit]\label{lem:orbit}
Starting from $\Btilde_0$ in \eqref{eq:B0extended}, perform the alternating
mutations $1,2,1,2,\ldots$. After $r\geq0$ mutations, the principal part is
$(-1)^rB(T_0)$ and the coefficient row is
\begin{equation}\label{eq:shearorbit}
 \ell_r=\begin{cases}
 (r+1,-r),&r\text{ even},\\
 (-r,r+1),&r\text{ odd}.
 \end{cases}
\end{equation}
Every matrix in this orbit is realized by the same one-component lamination
$L$ on the same annulus.
\end{lemma}
\begin{proof}
For $r=0$, the claim is \eqref{eq:B0extended}. Suppose first that $r$ is even.
The next mutation is at $1$. It negates the first coefficient $r+1$, and
changes the second from $-r$ to $-r+2(r+1)=r+2$. The top matrix changes sign.
This gives $(-(r+1),r+2)$, as required at the odd index $r+1$. If $r$ is odd,
the next mutation is at $2$, and the same calculation with the coordinates
interchanged gives $(r+2,-(r+1))$. Induction proves the formulas. The geometric
assertion follows by keeping $L$ fixed throughout and applying the
shear-coordinate mutation theorem at each flip.
\end{proof}

Fix an integer $d\geq1$. Choose the triangulation reached after $d-1$ mutations
as a new initial seed, and specialize its two mutable variables to $1$.
Continue forward along the alternating orbit, without further resetting any
variable. Write the successive specialized arc variables as
$P_2^{(d)},P_3^{(d)},\ldots$, and set $P_0^{(d)}=P_1^{(d)}=1$ for the two
initial variables in their exchange-chain order.

\begin{proposition}[Nonlinear recurrence]\label{prop:nonlinear}
The polynomials associated to this construction satisfy
\begin{equation}\label{eq:nonlinear}
 P_{n-1}^{(d)}P_{n+1}^{(d)}
   =(P_n^{(d)})^2+q^{d+n-1}\qquad(n\geq1).
\end{equation}
In particular $P_2^{(d)}=1+q^d$. All the terms arise from the same fixed curve
$L$, although changing $d$ changes the chosen initial triangulation.
\end{proposition}
\begin{proof}
Immediately before the $n$th forward exchange, the orbit index is
$r=d+n-2$. The column to be mutated has coefficient $r+1=d+n-1$ by
\eqref{eq:shearorbit}. Its nonzero mutable entry is $-2$: for even $r$ this is
the first column, and for odd $r$ it is the second. Thus the two exchange
monomials are $q^{d+n-1}$ and the square of the other mutable variable.
The exchanged-out variable is $P_{n-1}^{(d)}$ and the retained variable is
$P_n^{(d)}$. Equation~\eqref{eq:exchange} is exactly \eqref{eq:nonlinear}.
\end{proof}

The proof of the next section also shows directly that all these rationally
defined expressions are polynomials with nonnegative integer coefficients.
Consequently no general positivity theorem is needed for the enumeration
that follows.

\section{Linearization and positive enumeration}\label{sec:enumeration}

\subsection{An invariant that removes the division}

Set
\begin{equation}\label{eq:Kd}
 K_d(q)=1+q+q^d.
\end{equation}
Despite the quadratic form of \eqref{eq:nonlinear}, the sequence satisfies a
second-order linear recurrence with a fixed polynomial coefficient.

\begin{theorem}[Linear recurrence and generating functions]\label{thm:linear}
For every $d\geq1$,
\begin{equation}\label{eq:linear}
 P_0^{(d)}=P_1^{(d)}=1,\qquad
 P_{n+1}^{(d)}=K_d(q)P_n^{(d)}-qP_{n-1}^{(d)}\quad(n\geq1).
\end{equation}
The ordinary generating functions are
\begin{align}
 \sum_{n\geq0}P_n^{(d)}(q)t^n
   &=\frac{1-(q+q^d)t}{1-(1+q+q^d)t+qt^2},\label{eq:G}\\
 \sum_{N\geq0}P_{N+1}^{(d)}(q)t^N
   &=\frac{1-qt}{1-(1+q+q^d)t+qt^2}.\label{eq:H}
\end{align}
\end{theorem}
\begin{proof}
Temporarily work in the rational function field and abbreviate
$P_n=P_n^{(d)}$. Define
\[
 K_n=\frac{P_{n+1}+qP_{n-1}}{P_n}\qquad(n\geq1).
\]
For $n\geq2$, multiplication by $P_nP_{n-1}$ gives
\begin{align*}
 P_nP_{n-1}(K_n-K_{n-1})
 &=P_{n+1}P_{n-1}-P_n^2
    -q\bigl(P_nP_{n-2}-P_{n-1}^2\bigr)\\
 &=q^{d+n-1}-q\,q^{d+n-2}=0.
\end{align*}
Thus $K_n=K_1=P_2+q=1+q+q^d$. The rational functions in question are nonzero:
for positive $q$ the original exchange recursion gives positive values.
Rearranging $K_n=K_d$ proves \eqref{eq:linear}.

Conversely, define $P_n$ by \eqref{eq:linear} and put
$D_n=P_{n-1}P_{n+1}-P_n^2$. Direct substitution yields
\[
 D_n=qD_{n-1}\quad(n\geq2),\qquad D_1=P_2-1=q^d.
\]
Hence $D_n=q^{d+n-1}$, recovering the nonlinear identity. In particular the
linear recurrence is not merely a heuristic fit to initial terms.

Finally, multiplying the generating series by $1-K_dt+qt^2$ leaves only its
constant and linear coefficients. These are $1$ and $1-K_d=-(q+q^d)$,
respectively, giving \eqref{eq:G}. Subtracting its constant term and dividing
by $t$ gives \eqref{eq:H}.
\end{proof}

\subsection{A nonnegative transfer matrix}

The negative term in \eqref{eq:linear} conceals positivity, but a $2\times2$
transfer matrix makes it transparent. Define
\begin{equation}\label{eq:transfer}
 M_d(q)=\begin{pmatrix}1+q^d&q^d\\q&q\end{pmatrix}.
\end{equation}

\begin{proposition}[Positive walk model]\label{prop:walk}
For every $n\geq1$,
\begin{equation}\label{eq:matrixpower}
 P_n^{(d)}(q)=\bigl(M_d(q)^{n-1}\bigr)_{11}.
\end{equation}
Equivalently, $P_n^{(d)}$ counts length-$n-1$ walks that start and end at state
$A$ in the following weighted directed multigraph:
\begin{center}
\begin{tabular}{lll}
\toprule
Transition & Number of edges & Integer edge weights\\
\midrule
$A\to A$ & $2$ & $0$ and $d$\\
$A\to B$ & $1$ & $d$\\
$B\to A$ & $1$ & $1$\\
$B\to B$ & $1$ & $1$\\
\bottomrule
\end{tabular}
\end{center}
A walk of total edge weight $k$ contributes $q^k$. In particular,
$P_n^{(d)}\in\Z_{\geq0}[q]$.
\end{proposition}
\begin{proof}
The trace of $M_d$ is $K_d$ and its determinant is $q$. Its characteristic
identity is therefore $M_d^2-K_dM_d+qI=0$. The $(1,1)$ entries of its powers
satisfy \eqref{eq:linear}, with initial entries $1$ and $1+q^d$. These agree
with $P_1$ and $P_2$, proving \eqref{eq:matrixpower}. The entries of $M_d$
are the generating polynomials for the displayed transitions, so ordinary
matrix multiplication proves the walk interpretation. Each counted walk
has a nonnegative integer multiplicity.
\end{proof}

For $d\geq2$ and $n\geq2$, this also gives
\begin{equation}\label{eq:degree}
 \deg P_n^{(d)}=d(n-1),\qquad
 [q^0]P_n^{(d)}=[q^{d(n-1)}]P_n^{(d)}=1.
\end{equation}
Indeed, only the all-zero $A$-loop walk has weight zero. A walk of length
$n-1$ has weight at most $d(n-1)$; attaining this upper bound requires every
edge to have weight $d$. A walk returning from $B$ would use an edge of
weight $1<d$, so the only maximizing walk consists of the positive $A$-loop
at every step. For $d=1$ the upper-degree uniqueness argument does not apply,
which is why that case is excluded from the monicity assertion.

\subsection{Excursions and a positive finite formula}

A closed walk at $A$ decomposes uniquely into individual $A$-loops and
excursions from $A$ through $B$ back to $A$. If an excursion has $b\geq0$
intermediate $B$-loops, its length is $b+2$ and its weight is $d+b+1$.
Its generating series is thus
\[
 \sum_{b\geq0}t^{b+2}q^{d+b+1}
 =\frac{q^{d+1}t^2}{1-qt}.
\]
A closed walk is a sequence of these blocks and the two possible $A$-loops.
This gives a second derivation of \eqref{eq:H}:
\begin{equation}\label{eq:excursionGF}
 \sum_{N\geq0}P_{N+1}^{(d)}t^N
 =\frac{1}{1-(1+q^d)t-q^{d+1}t^2/(1-qt)}.
\end{equation}

\begin{theorem}[A subtraction-free finite expansion]\label{thm:finite}
For $N\geq0$,
\begin{align}\label{eq:positiveformula}
 P_{N+1}^{(d)}(q)
 &=(1+q^d)^N\\[-0.2em]
 &\quad+\sum_{j=1}^{\lfloor N/2\rfloor}\;
  \sum_{b=0}^{N-2j}
  \binom{N-j-b}{j}\binom{b+j-1}{j-1}
  q^{j(d+1)+b}(1+q^d)^{N-2j-b}.\notag
\end{align}
\end{theorem}
\begin{proof}
The first term counts walks without excursions. Otherwise let $j$ be the
number of excursions, let $b$ be the total number of internal $B$-loops,
and let $a=N-2j-b$ be the number of $A$-loops. There are
$\binom{a+j}{j}$ ways to position the $j$ ordered excursions among the $a$
individual $A$-loops. There are $\binom{b+j-1}{j-1}$ ways to distribute $b$
internal loops among the $j$ excursions. Their combined weight is
$j(d+1)+b$, and the $A$-loops contribute $(1+q^d)^a$. Summing gives
\eqref{eq:positiveformula}, since $a+j=N-j-b$.
\end{proof}

Thus these are ordinary weighted counting polynomials for a fixed finite
state system, in addition to being specializations of surface cluster
variables. This makes the recurrence and coefficient triangle accessible
without any cluster-algebra terminology.

\section{Infinitely many positive-coefficient obstructions}\label{sec:low}

Let $a_{n,k}^{(d)}=[q^k]P_n^{(d)}$. The beginning of each row can be read
directly from the walk model.

\begin{proposition}[Low-degree coefficients]\label{prop:low}
For $d\geq2$ and $n\geq1$,
\begin{align}
 a_{n,0}^{(d)}&=1,&
 a_{n,k}^{(d)}&=0\quad(1\leq k<d),\label{eq:low0}\\
 a_{n,d}^{(d)}&=n-1,\label{eq:lowd}\\
 a_{n,d+s}^{(d)}&=\max(n-1-s,0)
                 &&(1\leq s\leq d-1),\label{eq:lows}\\
 a_{n,2d}^{(d)}&=\binom{n-1}{2}+\max(n-d-1,0).
 \label{eq:low2d}
\end{align}
Here $\binom{0}{2}=\binom{1}{2}=0$.
\end{proposition}
\begin{proof}
Put $N=n-1$. The smallest positive $A$-loop weight is $d$ and the smallest
excursion weight is $d+1$. This proves \eqref{eq:low0}. A weight-$d$ walk
must contain one positive $A$-loop and otherwise zero loops, giving $N$
choices and \eqref{eq:lowd}.

For weight $d+s<2d$, there cannot be two positive $A$-loops, two excursions,
or one positive $A$-loop together with an excursion. The only possibility
is one excursion with $s-1$ internal $B$-loops. It occupies $s+1$ steps and
can be placed in $N-s$ positions if $N\geq s+1$, and in no positions
otherwise. This gives \eqref{eq:lows}.

At weight $2d$, there are exactly two possibilities. Two positive $A$-loops
can be chosen in $\binom N2$ ways. Alternatively, a single excursion with
$d-1$ internal $B$-loops occupies $d+1$ steps and has
$\max(N-d,0)$ possible positions. Their sum is \eqref{eq:low2d}.
\end{proof}

\begin{corollary}[Every $d\geq2$ supplies a positive triple]\label{cor:generaltriple}
At $n=d+1$, the coefficients of $P_n^{(d)}$ at consecutive degrees
$2d-2,2d-1,2d$ are
\begin{equation}\label{eq:generaltriple}
 \left(2,1,\binom d2\right).
\end{equation}
They violate log-concavity for every $d\geq2$. For $d\geq3$ they also form
a strict local valley.
\end{corollary}
\begin{proof}
For $d=2$, the first entry follows from \eqref{eq:lowd}; for $d\geq3$ it
follows from \eqref{eq:lows} with $s=d-2$. The remaining entries follow
from \eqref{eq:lows} and \eqref{eq:low2d}. The log-concavity inequality
would require $1\geq2\binom d2=d(d-1)$, which is false for every $d\geq2$.
For $d\geq3$, both neighbors of the middle entry exceed one.
\end{proof}

More is true without varying the initial seed at all.

\begin{proof}[Proof of \cref{thm:main}]
Use the fixed initial seed $T_1$, so $d=2$. The geometric realization and
recurrence follow from \cref{lem:realization,prop:nonlinear,thm:linear}.
For $n\geq3$, \cref{prop:low} gives the positive triple
\begin{equation}\label{eq:fixedtriple}
 \left(a_{n,2}^{(2)},a_{n,3}^{(2)},a_{n,4}^{(2)}\right)
 =\left(n-1,n-2,\binom{n-1}{2}+n-3\right).
\end{equation}
Write $m=n-3\geq0$. The neighbor product minus the middle square equals
\begin{align}\label{eq:defect}
 &(n-1)\left(\binom{n-1}{2}+n-3\right)-(n-2)^2\\
 &\hspace{2em}=\frac{m^3+5m^2+8m+2}{2}>0.\notag
\end{align}
Therefore log-concavity fails for every $n\geq3$. If $n\geq4$, the first
entry in \eqref{eq:fixedtriple} exceeds the second by one, while the third
exceeds the second by $\binom{n-1}{2}-1>0$. Thus the triple is a strict
valley. The expression for $P_4$ was calculated in \cref{sec:small}.
The degrees $2(n-1)$ strictly increase with $n$, so these polynomials are
indeed infinitely many distinct examples.
\end{proof}

For reference, the first coefficient rows for $d=2$ are shown below. Blank
entries are outside the degree range, not additional coefficients.
\begin{center}
\small
\setlength{\tabcolsep}{4.7pt}
\begin{tabular}{r|rrrrrrrrrrr|r}
\toprule
$n\backslash k$&0&1&2&3&4&5&6&7&8&9&10&$P_n(1)$\\
\midrule
0&1&&&&&&&&&&&1\\
1&1&&&&&&&&&&&1\\
2&1&0&1&&&&&&&&&2\\
3&1&0&2&1&1&&&&&&&5\\
4&1&0&3&2&4&2&1&&&&&13\\
5&1&0&4&3&8&7&7&3&1&&&34\\
6&1&0&5&4&13&14&20&16&11&4&1&89\\
\bottomrule
\end{tabular}
\end{center}

\section{Closed forms, Fibonacci numbers, and asymptotics}\label{sec:asymptotics}

\subsection{Exact algebraic form and row sums}

Fix $d\geq1$ and $q>0$, and define
\begin{equation}\label{eq:roots}
 \Delta_d(q)=\sqrt{K_d(q)^2-4q},\qquad
 \lambda_\pm(q)=\frac{K_d(q)\pm\Delta_d(q)}{2}.
\end{equation}
Since $K_d>1+q\geq2\sqrt q$, the roots are real, distinct, and positive.
The characteristic polynomial $x^2-K_dx+q$ is negative at both $x=1$ and
$x=q$. Therefore
\begin{equation}\label{eq:rootbounds}
 0<\lambda_-<\min(1,q)\leq\max(1,q)<\lambda_+.
\end{equation}

\begin{proposition}[Closed form]\label{prop:closed}
For every $n\geq0$ and $q>0$,
\begin{equation}\label{eq:closed}
 P_n^{(d)}(q)=
 \frac{(1-\lambda_-)\lambda_+^n+
             (\lambda_+-1)\lambda_-^n}{\lambda_+-\lambda_-}.
\end{equation}
At $q=1$, independently of $d$,
\begin{equation}\label{eq:fibonacci}
 P_0^{(d)}(1)=1,\qquad
 P_n^{(d)}(1)=F_{2n-1}\quad(n\geq1),
\end{equation}
where $F_0=0$, $F_1=1$, and $F_{r+1}=F_r+F_{r-1}$.
\end{proposition}
\begin{proof}
The right side of \eqref{eq:closed} satisfies the linear recurrence, because
both characteristic roots do. Its values at $n=0$ and $n=1$ are one, so
uniqueness proves the formula. At $q=1$, the recurrence becomes
$P_{n+1}=3P_n-P_{n-1}$, with $P_0=P_1=1$. The odd-indexed Fibonacci numbers
satisfy this recurrence and the required initial values for $n\geq1$.
Alternatively, substitute $\lambda_+=\varphi^2$ and
$\lambda_-=\varphi^{-2}$, where $\varphi=(1+\sqrt5)/2$, into
\eqref{eq:closed} and use Binet's formula.
\end{proof}

The row sums are exactly OEIS A001519 with its existing indexing:
\[
 1,1,2,5,13,34,89,233,610,1597,\ldots.
\]
That entry defines the sequence by the same recurrence and initial values,
and identifies it as an odd-indexed Fibonacci subsequence~\cite{OEIS}.
Thus no new OEIS identifier is being claimed for these totals. The family
provides a weight refinement of that established sequence.

There is a useful interpretation of the independence from $d$. At $q=1$,
\[
 M_d(1)=\begin{pmatrix}2&1\\1&1\end{pmatrix}
\]
is independent of $d$. The set of walks is the same for every parameter;
only the statistic assigned to each walk changes. Coefficient-shape phenomena
therefore change while the row sums do not.

\subsection{Exponential asymptotics at fixed positive \texorpdfstring{$q$}{q}}

Define
\[
 A_d(q)=\frac{1-\lambda_-(q)}{\Delta_d(q)},\qquad
 C_d(q)=\frac{\lambda_+(q)-1}{\Delta_d(q)}.
\]
Both are positive by \eqref{eq:rootbounds}. Equation~\eqref{eq:closed} gives
\begin{equation}\label{eq:asymp}
 P_n^{(d)}(q)=A_d(q)\lambda_+(q)^n
 \left(1+\frac{C_d(q)}{A_d(q)}
              \left(\frac{\lambda_-(q)}{\lambda_+(q)}\right)^n\right).
\end{equation}
This is an exact two-term expansion, and hence in particular an exponentially
accurate leading asymptotic. For fixed $d$, it is uniform when $q$ ranges over
any compact subset of $(0,\infty)$: the positive continuous quantities stay
bounded away from their forbidden limits, and the root ratio has a maximum
strictly less than one on that compact set.

At $q=1$, the exact expression is especially simple:
\begin{equation}\label{eq:fibasy}
 P_n^{(d)}(1)=\frac{\varphi^{2n-1}+\varphi^{-(2n-1)}}{\sqrt5}
 =\frac{\varphi^{2n-1}}{\sqrt5}
   \left(1+\varphi^{-4n+2}\right)\qquad(n\geq1).
\end{equation}

\subsection{A Gaussian limit despite persistent non-log-concavity}

Choose uniformly among all length-$n-1$ walks from $A$ to $A$, counting the
two $A$-loops as distinct edges. Let $W_n^{(d)}$ be the total weight. Its
probability generating function is
\begin{equation}\label{eq:pgf}
 \EE\bigl[q^{W_n^{(d)}}\bigr]
 =\frac{P_n^{(d)}(q)}{P_n^{(d)}(1)}.
\end{equation}

\begin{theorem}[Mean, variance, and central limit]\label{thm:clt}
For each fixed $d\geq1$, as $n\to\infty$,
\begin{align}
 \EE W_n^{(d)}&=\mu_dn+O(1),&
 \mu_d&=\frac12+\frac{2d-1}{2\sqrt5},\label{eq:mean}\\
 \Var(W_n^{(d)})&=\sigma_d^2n+O(1),&
 \sigma_d^2&=\frac{\sqrt5(2d^2-2d+3)}{25}.
 \label{eq:variance}
\end{align}
Furthermore,
\begin{equation}\label{eq:clt}
 \frac{W_n^{(d)}-\mu_dn}{\sqrt{\sigma_d^2n}}
 \ \Longrightarrow\ N(0,1).
\end{equation}
\end{theorem}
\begin{proof}
The discriminant in \eqref{eq:roots} is nonzero at $q=1$. Choose analytic
branches of the roots in a complex neighborhood of that point. Their moduli
remain separated in a sufficiently small neighborhood, and $A_d$ is nonzero
there. If $s$ is sufficiently close to zero, put
$H(s)=\log\lambda_+(e^s)$, choosing the analytic logarithm through its positive
real value at zero. From \eqref{eq:asymp}, uniformly in a smaller neighborhood,
\begin{equation}\label{eq:quasipower}
 \log\frac{P_n^{(d)}(e^s)}{P_n^{(d)}(1)}
 =n\bigl(H(s)-H(0)\bigr)
  +\log\frac{A_d(e^s)}{A_d(1)}+O(\rho^n),
 \qquad 0<\rho<1.
\end{equation}
Analyticity and a smaller disk also permit differentiation of the exponentially
small remainder. Thus the first two cumulants have leading coefficients
$H'(0)$ and $H''(0)$.

We compute them explicitly rather than leaving them as derivatives of a root.
The characteristic equation becomes
\[
 e^{H(s)}+e^{s-H(s)}=1+e^s+e^{ds}.
\]
Put $L=(3+\sqrt5)/2$, $l=(3-\sqrt5)/2$, $\mu=H'(0)$, and $\nu=H''(0)$.
The first and second derivatives at zero give
\begin{align*}
 L\mu+l(1-\mu)&=1+d,\\
 L(\mu^2+\nu)+l\bigl((1-\mu)^2-\nu\bigr)&=1+d^2.
\end{align*}
Since $L-l=\sqrt5$, the first equation yields \eqref{eq:mean}. Substituting
that value into the second yields
\[
 \nu=\frac{1+d^2-L\mu^2-l(1-\mu)^2}{\sqrt5}
     =\frac{\sqrt5(2d^2-2d+3)}{25}>0.
\]
This proves the mean and variance formulas.

For the limit law, expand
$H(s)-H(0)=\mu_ds+\tfrac12\sigma_d^2s^2+O(s^3)$ in
\eqref{eq:quasipower}, set $s=iu/\sqrt{\sigma_d^2n}$ for fixed real $u$,
and subtract the centering term $iu\mu_dn/\sqrt{\sigma_d^2n}$. The logarithm
of the resulting characteristic function tends to $-u^2/2$; the analytic
prefactor contributes $O(n^{-1/2})$ and the cubic remainder does likewise.
The continuity theorem for characteristic functions proves
\eqref{eq:clt}.
\end{proof}

For the fixed-seed counterexample family $d=2$, the constants are
\[
 \mu_2=\frac12+\frac{3\sqrt5}{10},\qquad
 \sigma_2^2=\frac{7\sqrt5}{25}.
\]
There is no contradiction between this Gaussian limit and the coefficient
valleys proved earlier. The central limit theorem describes weight on the
scale $n$ with fluctuations of order $\sqrt n$. Our persistent obstruction
occurs at the fixed degrees $2,3,4$, far into the lower tail. Those
coefficients grow only polynomially in $n$, whereas the row sum grows
exponentially. A Gaussian weak limit neither asserts local log-concavity nor
controls every fixed tail coefficient.

\section{Interpretation, verification, and remaining questions}\label{sec:scope}

\subsection{What the counterexample does and does not show}

The decisive distinction is between \emph{one connected curve} and
\emph{at most one contributing crossing at each mutable arc}. The former does
not imply the latter on an annulus. The coefficient row $(-1,2)$ is obtained
from one curve by one flip, and \eqref{eq:shearorbit} exhibits unbounded shear
coordinates with no change in the number of curve components. The resulting
exchange relation contains $q^2$, then $q^3$, then $q^4$.

This identifies a concrete obstruction to extending arguments based only on
unit coefficient shifts to arbitrary unpunctured surfaces. It is unnecessary
to diagnose every line of an existing proof: the explicit example already
settles the general formulation. The mechanism, however, explains why tests
confined to polygons could miss it.

The natural residual question is the polygon-restricted formulation: does
log-concavity hold for every single lamination in an unpunctured polygon,
with the same geometric-coefficient specialization? This note does not answer
that question. Nor does it characterize exactly which triples $(S,T,L)$
produce log-concave polynomials. A further possible restriction is to bound
the shear-coordinate multiplicities relevant to an expansion, but the
sufficiency of such a condition would require a new proof; it is not asserted
here.

The annular construction is deliberately narrow enough to be audited. It
uses a rank-two exchange matrix, a single unit-weight bridge, legal flips,
and ordinary integer arithmetic. It does not rely on a large experimental
search, a floating-point inequality, or a numerical guess about a lamination's
connectedness.

\subsection{Exact computational verification}

The supplied \texttt{verify.py} uses the Python standard library and represents
polynomials as tuples of integers in increasing degree. Polynomial division
checks both integrality of the quotient and vanishing of the remainder.
It implements the full extended-matrix mutation rule separately from the
linear and nonlinear recurrences and separately from the walk update.
The following checks were actually executed:
\begin{center}
\small
\begin{tabular}{p{0.64\linewidth}r}
\toprule
Check & Number of cases\\
\midrule
Shear-orbit matrices, with both mutation involutions & 201\\
Linear/walk agreement, positivity, and Fibonacci row sums & 972\\
Nonlinear recurrence and actual seed mutations & 396\\
Positive finite expansion against the recurrence & 156\\
Fixed-seed log-concavity failures, $3\leq n\leq160$ & 158\\
Fixed-seed strict valleys, $4\leq n\leq160$ & 157\\
\bottomrule
\end{tabular}
\end{center}
The 972 cases are $1\leq d\leq12$ and $0\leq n\leq80$. The 396
nonlinear/mutation cases use the same $d$ range and $0\leq n\leq32$.
The finite-expansion comparison uses $0\leq n\leq12$. All comparisons are
exact. The detailed low-degree formulas, degree, monicity, and the general
triple \eqref{eq:generaltriple} are also checked in their stated ranges.

These computations do not machine-formalize the geometric theorem or certify
all $n$. The article supplies the geometric reasoning and the induction and
counting proofs that make the conclusions unconditional for the stated
infinite families. No Lean formalization or external computer-algebra
certification is claimed.

\subsection{Reproducibility and research status}

To rerun the checks and regenerate the data, use Python~3.10 or later:
\begin{lstlisting}
python verify.py --write-data
\end{lstlisting}
The program writes a matrix-and-polynomial certificate, a coefficient table
for $d=2$ through $n=40$, Fibonacci row sums through $n=100$, and a JSON
verification report. The archive also includes the LaTeX source, this PDF,
a source-audit note, and build instructions. The final proof does not depend
on reproducing a random search.

As of the research date, the arXiv record inspected still presented the target
as a conjecture. Targeted searches by title and identifier did not locate a
correction addressing this annular example. That is not an exhaustive
priority determination. Before preparing a journal submission, the appropriate
next verification is to compare the final publisher version and seek expert
review of the interpretation of the geometric coefficients. The counterexample
and all algebraic claims made here are stated explicitly enough for that
review to be short and decisive.

\appendix
\section{A minimal independent check of the fixed-seed family}\label{app:minimal}

The following small program uses only coefficientwise recurrence updates and
is included as \texttt{minimal\_check.py}. It is not the main verifier, since
it does not implement the surface matrix. Its purpose is to make the central
coefficient claim reproducible with very little code.
\begin{lstlisting}[language=Python]
def add(a, b):
    c = [0] * max(len(a), len(b))
    for i, v in enumerate(a):
        c[i] += v
    for i, v in enumerate(b):
        c[i] += v
    while len(c) > 1 and c[-1] == 0:
        c.pop()
    return c

def shift(a, k):
    return [0] * k + a

P = [[1], [1]]
for n in range(1, 160):
    next_p = add(add(P[n], shift(P[n], 1)), shift(P[n], 2))
    next_p = add(next_p, [-v for v in shift(P[n-1], 1)])
    P.append(next_p)

assert P[4] == [1, 0, 3, 2, 4, 2, 1]
assert all(p[3]**2 < p[2]*p[4] for p in P[3:])
assert all(p[2] > p[3] < p[4] for p in P[4:])
print("PASS: exact fixed-seed checks through n=160")
\end{lstlisting}

\section{Certificate at a glance}\label{app:certificate}

The essential proof data can be summarized without any asymptotic or
generating-function machinery:
\begin{center}
\begin{tabular}{p{0.29\linewidth}p{0.62\linewidth}}
\toprule
Surface & Annulus, one marked point per boundary, no punctures.\\
Reference triangulation & Two bridges forming two triangles; principal
matrix $\bigl(\begin{smallmatrix}0&2\\-2&0\end{smallmatrix}\bigr)$.\\
Lamination & One elementary curve along the first bridge, with weight one;
reference shear row $(1,0)$.\\
Geometric certification & Flip reference arc $1$, keeping the same curve;
new coefficient row $(-1,2)$.\\
Chosen initial variables & The two mutable variables $a,b$ of this new seed;
set $a=b=1$ only when forming $c$-polynomials.\\
Subsequent mutations & $2,1,2$.\\
Exchange equations & $bv=a^2+q^2$, $aw=v^2+q^3$, $vz=w^2+q^4$.\\
Resulting polynomial & $1+3q^2+2q^3+4q^4+2q^5+q^6$.\\
Violation & $2^2<3\cdot4$, using positive coefficients at degrees $2,3,4$.\\
\bottomrule
\end{tabular}
\end{center}

\begin{thebibliography}{9}

\bibitem{KLL}
W.~Kang, K.~J.~Lee, and E.~Lim,
\emph{Unimodality and Cluster Algebras from Surfaces},
arXiv:2508.04396v3, 11 February 2026.
Target statements: Definition~20, Theorem~14, and Conjecture~1 in Section~5.
\url{https://arxiv.org/abs/2508.04396v3};
\url{https://arxiv.org/html/2508.04396v3}.
The arXiv record lists the related journal publication
\emph{European Journal of Combinatorics} \textbf{136} (2026), 104392,
\url{https://doi.org/10.1016/j.ejc.2026.104392}.
Version-specific text and metadata consulted 19 September 2026.

\bibitem{FT}
S.~Fomin and D.~Thurston,
\emph{Cluster algebras and triangulated surfaces. Part II: Lambda lengths},
\emph{Memoirs of the American Mathematical Society} \textbf{255} (2018),
no.~1223, v+97 pp.
\url{https://arxiv.org/abs/1210.5569}.
Especially Definitions~12.5 and~17.2, Theorems~12.6 and~13.5,
and Proposition~17.3.

\bibitem{FST}
S.~Fomin, M.~Shapiro, and D.~Thurston,
\emph{Cluster algebras and triangulated surfaces. I. Cluster complexes},
\emph{Acta Mathematica} \textbf{201} (2008), 83--146.
\url{https://arxiv.org/abs/math/0608367}.

\bibitem{FZ}
S.~Fomin and A.~Zelevinsky,
\emph{Cluster algebras. I. Foundations},
\emph{Journal of the American Mathematical Society} \textbf{15} (2002),
497--529.
\url{https://arxiv.org/abs/math/0104151}.

\bibitem{OEIS}
The OEIS Foundation Inc.,
\emph{The On-Line Encyclopedia of Integer Sequences}, entry A001519,
$a(0)=a(1)=1$, $a(n)=3a(n-1)-a(n-2)$.
\url{https://oeis.org/A001519}.
Consulted 19 September 2026.

\end{thebibliography}
\end{document}
